\section{Local weight adjustment with bounded participation}
\label{sec:adjustment}

The coloring argument requires local changes of weights that offset the
contribution of triangles.  Each change must also involve few entries
from any one list.  We formulate this requirement for a partition of a
weighted graph, with constants independent of the graph and of the size
of its parts.

Fix an integer $r\ge4$.  Let $H$ be a finite $K_r$-free graph with a
partition $\mathcal D$ of its vertex set such that
\[
 |N_H(i)\cap D|\le1
 \qquad(i\in V(H),\ D\in\mathcal D).
\]
Give each vertex $i$ a positive weight $w_i$, and write
\begin{equation}
\label{eq:weighted-local-masses}
 s_i=\sum_{u\in N_H(i)}w_u,\qquad
 d_i=w_i s_i,\qquad
 t_i=w_i\sum_{u\in N_H(i)}w_u
                 \sum_{z\in N_H(u)\cap N_H(i)}w_z.
\end{equation}
Thus $d_i$ is the mass of edges incident to $i$, and $t_i$ is twice the
mass of triangles containing $i$.  Suppose that $s_i\le\Lambda$ for
every $i$, where $\Lambda\ge1$.  Define
\[
 f(a)=a\log a-a+1\quad(a>0),\qquad f(0)=1,
\]
and put
\begin{equation}
\label{eq:adjustment-parameters}
 \ell_0=\log(8\Lambda),\qquad
 R_0=\lceil4\ell_0\rceil+10,\qquad
 U=\exp(A\ell_0^2),\qquad
 P_0=\left\lceil\exp(40\ell_0^2)\right\rceil,
\end{equation}
where $A=A(r)>0$ will be chosen below.

A \emph{multiplier operation} consists of a nonempty listed support
$J\subseteq V(H)$, contained in a ball of radius $R_0$ in $H$ and
satisfying $|J\cap D|\le P_0$ for every $D\in\mathcal D$, together with
multipliers $m_i\in[0,U]$ for $i\in J$.  Extend $m_i=1$ off $J$, and
associate to the operation the vector
\begin{equation}
\label{eq:multiplier-vector}
 g_i=w_i f(m_i)+\sum_{u\in N_H(i)}
                       w_iw_u(m_i-1)(m_u-1).
\end{equation}
Multipliers equal to $1$ are permitted on $J$.  A \emph{drop operation}
at $i$ has listed support $J=\{i\}$ and vector
\begin{equation}
\label{eq:drop-vector}
 g_i=\Lambda w_i-d_i,\qquad
 g_u=-w_iw_u\quad(u\in N_H(i)),\qquad
 g_u=0\quad(u\notin N_H(i)\cup\{i\}).
\end{equation}
The listed support records the vertices on which the operation acts; a
drop vector can have nonzero coordinates outside that support.  It charges
$\Lambda w_i$ for losing weight at $i$ and subtracts its incident edge
masses.

The two terms of \eqref{eq:multiplier-vector} have a first-order
interpretation.  For fixed multipliers $m_i$ and small $\rho>0$, apply
them on a common event of probability $\rho$, and use
$(1-\rho m_i)/(1-\rho)$ otherwise.  Each weight then has unchanged
expectation.  The expected change of the entropy term
$w_i\log w_i-w_i$, with $0\log0=0$, is
$\rho w_i f(m_i)+O(\rho^2)$, while the
expected change on an edge $iu$ is
$\rho w_iw_u(m_i-1)(m_u-1)+O(\rho^2)$.  In particular, events
of probability $\rho w_z$ that delete the weights on $N_H(z)$,
with this compensation of their mean loss, contribute $\rho d_i$ to
the entropy change at $i$ and $\rho t_i$ to its incident edge change,
to first order.  The following lemma supplies operations that offset
the latter contribution simultaneously at all vertices.

\begin{lemma}[Local weight adjustment]
\label{lem:adjustment}
For every integer $r\ge4$ there are constants $A=A(r)>0$,
$B=B(r)\ge1$, and $\Lambda_r\ge1$ with the following property.
Let $\Lambda\ge\Lambda_r$, and let $H$ be a finite $K_r$-free graph
whose vertex set has a partition $\mathcal D$ such that each vertex has
at most one neighbor in each part.  For any positive vertex weights
with $s_i\le\Lambda$, there is a finite family of the multiplier and
drop operations defined above, with vectors $g^{(j)}$, listed supports
$J_j$, and rates $a_j\ge0$, such that
\begin{equation}
\label{eq:adjustment}
 t_i+\sum_j a_jg_i^{(j)}\le B d_i,
 \qquad
 \sum_{j:\,i\in J_j}a_j\le\Lambda^4
 \qquad(i\in V(H)).
\end{equation}
The constants depend only on $r$.
\end{lemma}

Taking every part to be a singleton gives the same weight adjustment
conclusion for an arbitrary positively weighted $K_r$-free graph, with
no partition specified: the neighbor and support conditions on parts
then hold automatically.

The proof occupies this section and Section~\ref{sec:triangles}.
Nonnegative prices first reduce the simultaneous vertex inequalities to
a scalar triangle estimate.  We then derive a local edge estimate from
the same prices.  Section~\ref{sec:triangles} uses that edge estimate to
bound the triangles and finish the proof.  If $H$ has no vertices, the
lemma holds with no operations, so we assume otherwise throughout the
reduction.

\subsection{Dual prices and dyadic levels}

A bound on total triangle mass would not guarantee the individual
vertex inequalities in \eqref{eq:adjustment}.  To retain those
constraints, assign prices $x_i\ge0$ to the first inequalities in
\eqref{eq:adjustment}, and prices $w_i y_i\ge0$ to the incidence
budgets in the second inequalities.  It suffices to prove that
\begin{equation}
\label{eq:dual-test}
 \sum_i x_i g_i+\sum_{i\in J}w_i y_i\ge0
 \qquad\text{for every allowed operation $(g,J)$}
\end{equation}
implies
\begin{equation}
\label{eq:dual-target}
 \sum_i x_i t_i\le B\sum_i x_i d_i
                         +\Lambda^4\sum_i w_i y_i.
\end{equation}
The support price in \eqref{eq:dual-test} includes every vertex of the
listed support, even when its multiplier equals $1$.

To justify the reduction, let $n=|V(H)|$ and let $\mathcal C$ be the
cone in $\mathbb R^{2n}$ generated by the operation columns
$(g,\mathbf1_J)$ and the nonnegative orthant.  These columns form a
compact set: there are finitely many possible supports, each multiplier
cube is compact, and $f$ is continuous on $[0,U]$.  Moreover, the sum
of the last $n$ coordinates of every column is at least $1$.  In any
convergent sequence of elements of $\mathcal C$, the last-coordinate
sums therefore bound the total coefficients of the operation columns.
Their combined contributions belong to a fixed multiple of the compact
convex hull of the columns and zero.  A subsequence of those
contributions converges; the residual vectors then converge in the
nonnegative orthant.  This proves that $\mathcal C$ is closed.

Membership of $(Bd-t,\Lambda^4\mathbf1)$ in $\mathcal C$ is exactly
the existence of the finite family in \eqref{eq:adjustment}.  If this
target were outside the cone, separation would give a linear functional
nonnegative on $\mathcal C$ and negative at the target.  Since
$\mathcal C$ contains the nonnegative orthant, the coefficients of
that functional are nonnegative.  Because $w_i>0$, they may be written
as $(x_i,w_i y_i)_i$.  They would satisfy \eqref{eq:dual-test} and
violate \eqref{eq:dual-target}, proving the reduction.

Fix nonnegative prices satisfying \eqref{eq:dual-test}.  Applying that
inequality to the drop at $i$ and dividing by $w_i$ gives
\begin{equation}
\label{eq:drop-test}
 \sum_{u\in N_H(i)}x_u w_u
             \le(\Lambda-s_i)x_i+y_i\le\Lambda x_i+y_i.
\end{equation}
Call a vertex $i$ \emph{good} if $x_i>y_i$, and \emph{bad} otherwise.
A triangle $\{i,u,z\}$ contributes
$2(x_i+x_u+x_z)w_iw_uw_z$ to $\sum_v x_vt_v$.  For a fixed bad
vertex $i$, the sum of these contributions over triangles containing
$i$ is at most
\[
 2w_i\left(x_i s_i^2+
               s_i\sum_{u\in N_H(i)}x_u w_u\right)
 \le O(\Lambda^2)w_i y_i,
\]
by \eqref{eq:drop-test} and $s_i\le\Lambda$.  Summing over bad
vertices, with possible overcounting, yields
\begin{equation}
\label{eq:bad-triangles}
 \sum_{\substack{\{i,u,z\}\text{ triangle of }H\\
                 \{i,u,z\}\text{ contains a bad vertex}}}
     2(x_i+x_u+x_z)w_iw_uw_z
 \le O(\Lambda^2)\sum_i w_i y_i.
\end{equation}
The remaining triangles have only good vertices.  Their edge budget is
\begin{equation}
\label{eq:price-mass}
 M_x=\sum_i x_i d_i
     =\sum_{iu\in E(H)}(x_i+x_u)w_iw_u,
\end{equation}
where the edge sum is unordered.

Good vertices have positive prices.  Consider all dyadic numbers
$b=2^j$, $j\in\mathbb Z$, such that $[b,2b)$ contains the price of a
good vertex.  There are only finitely many such levels.  At level $b$,
call a good vertex low if $x_i<b$, middle if $b\le x_i<2b$, and high
if $x_i\ge2b$.  Let $H_b$ be the graph on good vertices obtained from
$H$ by deleting edges with two low endpoints or two high endpoints.
Rescale the weights to
\begin{equation}
\label{eq:rescaled-weights}
 q_i=
 \begin{cases}
  w_i x_i/b,&i\text{ high},\\
  w_i,&i\text{ low or middle}.
 \end{cases}
\end{equation}
The weights $q$ and all three classifications refer to the level under
consideration.

A triangle of good vertices survives at exactly one level: the level
whose interval contains its median price.  At that level it has at most
one low and at most one high vertex, and hence at least one middle
vertex.  At every other level, two vertices are low or two are high,
and the edge between them is deleted.  The rescaling pays for the
largest price even when it is much larger than the median: if
$x_i<b\le x_u<2b\le x_z$, then
\[
 bq_iq_uq_z=x_z w_iw_uw_z.
\]
If a retained triangle has no high vertex, all three prices are less than $2b$.
If it has a high vertex, its $q$-weight incorporates that largest price
divided by $b$.  In either case its contribution satisfies
\[
 2(x_i+x_u+x_z)w_iw_uw_z\le12bq_iq_uq_z.
\]
Define
\begin{equation}
\label{eq:dyadic-triangle-mass}
 T_0=\sum_b b
       \sum_{\{i,u,z\}\text{ triangle of }H_b}q_iq_uq_z.
\end{equation}
We have proved
\begin{equation}
\label{eq:triangle-reduction}
 \sum_i x_i t_i\le12T_0+O(\Lambda^2)\sum_i w_i y_i.
\end{equation}
Consequently, a bound $T_0=O_r(M_x)$ will imply
\eqref{eq:dual-target} once $B$ and the lower threshold on $\Lambda$
are chosen.

To obtain that bound, we express the dual prices in the rescaled
weights.  For good vertices and original edges between them, set
\begin{equation}
\label{eq:rescaled-coefficients}
 \alpha_i=\frac{x_iw_i}{bq_i},\qquad
 \beta_i=\frac{y_iw_i}{bq_i},\qquad
 \kappa_{iu}=\frac{(x_i+x_u)w_iw_u}{bq_iq_u}.
\end{equation}
They satisfy
\begin{equation}
\label{eq:coefficient-bounds}
 0\le\beta_i<\alpha_i\le2,\qquad
 0<\kappa_{iu}\le4,\qquad
 \kappa_{iu}\ge1\quad(iu\in E(H_b)).
\end{equation}
Indeed, $\alpha_i=1$ at high vertices and $\alpha_i=x_i/b$ at the
others.  If exactly one endpoint of an edge is high, then
$\kappa_{iu}$ is $1$ plus the ratio of the other price to the high
price.  If both endpoints are high it is $b/x_i+b/x_u$, and if neither
is high it is $(x_i+x_u)/b$.  A retained edge of the last type has a
middle endpoint, proving its lower bound.

For a multiplier operation supported on good vertices, division of
\eqref{eq:dual-test} by $b$ gives
\begin{equation}
\label{eq:normalized-dual}
 \begin{split}
 0\le{}&\sum_{i\text{ good}}\alpha_i q_i f(m_i)
   +\sum_{\substack{iu\in E(H)\\i,u\text{ good}}}
        \kappa_{iu}q_iq_u(m_i-1)(m_u-1)\\
   &+\sum_{i\in J}\beta_iq_i.
 \end{split}
\end{equation}
The edge sum here still uses the original edges of $H$, including
edges absent from $H_b$.

The rescaling also preserves a uniform neighborhood bound:
\begin{equation}
\label{eq:level-neighborhood}
 \sum_{u\in N_{H_b}(i)}q_u\le e^{\ell_0}=8\Lambda.
\end{equation}
For a high vertex, all retained neighbors have their original weights,
so the sum is at most $\Lambda$.  For a non-high vertex $i$, the high
neighbors have total $q$-weight at most
\[
 \frac1b\sum_{u\in N_H(i)}x_uw_u
 \le\frac{\Lambda x_i+y_i}{b}<2(\Lambda+1),
\]
by \eqref{eq:drop-test} and goodness.  The other neighbors contribute
at most $\Lambda$, proving \eqref{eq:level-neighborhood}.

\subsection{Sparse multipliers and local edge mass}

The normalized dual inequality controls the mass of edges between two
nearby sets.  To see this, we will increase the weights on one set and
set the weights on the other to zero.  The cross edges then have
negative quadratic contribution.  The following construction keeps the
positive contribution inside the first set small.  Its recursive passage
to neighborhoods excluding a smaller clique is the sparse-subgraph
device of Ajtai, Erd\H{o}s, Koml\'os, and
Szemer\'edi~\cite[Section~3]{AEKS1981}.  The mean-one weighted
construction and the resulting cross estimate adapt
\cite[Lemmas~2.2--2.3]{OpenAIIndependence}; we give the proofs with the
locality and participation conditions needed here.

\begin{lemma}[Sparse multipliers]
\label{lem:sparse-multipliers}
For every integer $r\ge2$ there is a constant $D_r\ge1$ such that the
following holds.  Let $F$ be a finite $K_r$-free graph with positive
vertex weights $q_i$ of total mass $s$.  For every
$0<\epsilon\le1/2$, there are jointly distributed random multipliers
$n_i$ satisfying
\begin{equation}
\label{eq:sparse-multipliers}
 0\le n_i\le D_r\epsilon^{-(r-2)},\qquad
 \EE n_i=1,\qquad
 \EE\sum_{iu\in E(F)}q_iq_u n_i n_u\le\epsilon s^2,
\end{equation}
where the edge sum is unordered.
\end{lemma}

\begin{proof}
The assertion is vacuous if $F$ has no vertices.  Otherwise $s>0$,
and we argue by induction on $r$.  For $r=2$, the graph has no edges
and $n_i=1$ suffices.

Keep the initial total mass $s$ fixed.  Successively remove the open
neighborhood of a vertex in the current remainder whenever that
neighborhood has mass greater than $\epsilon s/4$.  The removed sets
are disjoint, each induces a $K_{r-1}$-free graph, and their masses
$s_j$ satisfy $s_j>\epsilon s/4$.  In the final remainder, every
neighborhood has mass at most $\epsilon s/4$, so its unordered edge
mass is at most $\epsilon s^2/8$.  If no set was removed, the constant
multipliers $1$ already prove the assertion.

Otherwise, select the final remainder with probability $1/2$, or
removed set $j$ with probability
\[
 p_j=\frac{s_j}{2\sum_l s_l}>\frac\epsilon8.
\]
Set the multipliers to zero off the selected set.  On the selected
remainder, use the constant multiplier $2$.  On a selected removed set
$j$, use the inductively supplied multipliers with parameter
$\epsilon/4$, divided by $p_j$.  Every vertex has mean multiplier
$1$.  On a removed set, the multiplier bound follows from
\[
 \frac{D_{r-1}(\epsilon/4)^{-(r-3)}}{p_j}
 \le 8\cdot4^{r-3}D_{r-1}\epsilon^{-(r-2)}.
\]
Thus a constant $D_r$ depending only on $r$ bounds all multipliers as
required.

Edges between different sets contribute zero.  The remainder
contributes at most $\epsilon s^2/4$ in expectation, while the removed
sets together contribute at most
\[
 \sum_j\frac{(\epsilon/4)s_j^2}{p_j}
 =\frac\epsilon2\left(\sum_j s_j\right)^2
 \le\frac\epsilon2s^2.
\]
Their sum is at most $\epsilon s^2$, completing the induction.
\end{proof}

For sets $S,D$ in a graph with weights $q$, write
\[
 q(S)=\sum_{i\in S}q_i,\qquad
 e_q(S,D)=\sum_{i\in S}\sum_{u\in D\cap N(i)}q_iq_u.
\]
This is a directed edge sum: an edge with both endpoints in $S\cap D$
is counted twice.  We specify the graph by a subscript when needed.

\begin{lemma}[Local cross estimate]
\label{lem:local-cross}
Assume the prices satisfy \eqref{eq:dual-test}, and fix a level $b$.
Let $S,D\subseteq V(H_b)$, where $S\cup D$ is contained in a ball of
radius $R_0$ in $H$ and meets every part of $\mathcal D$ in at most
$P_0$ vertices.  Suppose that
\[
 32\le h\le20\ell_0^2,\qquad q(S),q(D)\le e^h,
 \qquad 0\le\eta\le2,
\]
and that $\alpha_i,\beta_i\le\eta$ for $i\in S$.  If $A$ in
\eqref{eq:adjustment-parameters} is sufficiently large depending only
on $r$, then
\begin{equation}
\label{eq:local-cross}
 e_{q,H_b}(S,D)\le C_1\bigl(\eta hq(S)+e^{-3h}\bigr),
\end{equation}
where $C_1=C_1(r)$.
\end{lemma}

\begin{proof}
First suppose that $S,D$ are disjoint and nonempty.  Apply
Lemma~\ref{lem:sparse-multipliers} to the original induced graph
$H[S]$, with its $q$-weights and $\epsilon=e^{-16h}$.  Set
$k'=e^{8h}$, and take the multiplier operation with listed support
$S\cup D$ and
\[
 m_i=1+k'n_i\quad(i\in S),\qquad m_i=0\quad(i\in D).
\]
The sparse multiplier bound implies
$\log(1+k'n_i)\le C'_r h$.  Since $h\le20\ell_0^2$, choosing
$A\ge20C'_r$ makes every realization an allowed operation; the ball
and participation conditions follow from the hypotheses on $S\cup D$.

Write $s=q(S)$ and $d=q(D)$, and take expectations in
\eqref{eq:normalized-dual}.  On $S$, the expected entropy term is at
most $C'_r\eta hk's$, because $f(1+a)\le a\log(1+a)$ for $a\ge0$
and $\EE n_i=1$.  The support price on $S$ is at most $\eta s$.
On $D$, the entropy and support terms together are at most $4d$,
using $f(0)=1$ and \eqref{eq:coefficient-bounds}.

The expected quadratic edge term is at most
\begin{equation}
\label{eq:cross-quadratic-cost}
 4(k')^2\epsilon s^2+2d^2-k'e_{q,H_b}(S,D).
\end{equation}
Inside $S$, the bound $\kappa_{iu}\le4$ and
\eqref{eq:sparse-multipliers} give the first term, including for
original edges absent from $H_b$.  Inside $D$, the unordered edge mass
is at most $d^2/2$, giving the second term.  On retained edges between
$S$ and $D$, the coefficient is at least $1$ and the expected product
of deviations is $-k'$.  Deleted cross edges have nonpositive
contributions and can be omitted in an upper bound.  Edges leaving
$S\cup D$ contribute zero.

Since the tested inequality holds for every realization, these bounds
give
\[
 e_{q,H_b}(S,D)
 \le C'_r\eta hs+
       \frac{\eta s+4d+4(k')^2\epsilon s^2+2d^2}{k'}
 \le C'_r\eta hs+O(e^{-3h}).
\]
For the last inequality, $s,d\le e^h$, $\eta\le2$,
$k'=e^{8h}$, and $(k')^2\epsilon=1$ make the numerator after
division at most $6e^{-7h}+6e^{-6h}$.  If either set is empty, the
bound holds immediately.

For possibly overlapping sets, assign every vertex independently to a
left or right class with equal probabilities.  Apply the disjoint
estimate to the part of $S$ on the left and the part of $D$ on the
right.  Each directed edge of $e_{q,H_b}(S,D)$ survives with
probability $1/4$, whereas the expected mass of the left subset of
$S$ is $q(S)/2$.  Taking expectations and increasing $C_1$ proves
\eqref{eq:local-cross} in general.
\end{proof}

The triangle argument will replace vertices by copies and delete some
edges.  The next observation ensures that the local estimate remains
available after these changes.  A homomorphism $p:F\to H_b$ is
\emph{injective on directed edges} if the map
$(u,v)\mapsto(p(u),p(v))$ is injective on all ordered adjacent pairs of
$F$.  Vertices may have several preimages, but each original directed
edge has at most one preimage.

\begin{corollary}[Transfer under vertex splitting]
\label{lem:local-cross-transfer}
Let $F$ be a graph with a homomorphism $p:F\to H_b$ that is injective
on directed edges.  Give every vertex of $F$ the weight, prices,
coefficients, and part of its image.  Then every vertex of $F$ has at
most one neighbor in each inherited part.  Moreover,
\eqref{eq:local-cross} holds in $F$ for any $S,D\subseteq V(F)$
satisfying the mass and coefficient hypotheses of
Lemma~\ref{lem:local-cross}, provided $S\cup D$ lies in a ball of
radius $R_0$ in $F$ and meets each inherited part in at most $P_0$
vertices.
\end{corollary}

\begin{proof}
Directed-edge injectivity implies
\[
 e_{q,F}(S,D)\le e_{q,H_b}(p(S),p(D)),\qquad
 q(p(S))\le q(S),\qquad q(p(D))\le q(D).
\]
The projected union lies in a ball of radius $R_0$ in $H$ and meets
each original part in at most $P_0$ vertices.  The coefficient bounds
also pass to the images.  Lemma~\ref{lem:local-cross}, which permits
overlap of the two sets, therefore proves the asserted estimate.
Finally, distinct neighbors of a vertex of $F$ have distinct images:
otherwise two directed edges would have the same image.  Since their
images lie in the neighborhood of one vertex of $H_b$, at most one
belongs to any given part.
\end{proof}

We have derived the local edge estimate solely from the dual condition
and have shown that it survives the graph modifications needed below.
The remaining assertion is the following triangle bound.

\begin{proposition}[Triangle bound]
\label{prop:triangle-bound}
For the graphs $H_b$, weights $q$, and prices constructed above from
\eqref{eq:dual-test}, with $A$ chosen as in
Lemma~\ref{lem:local-cross}, the quantities in
\eqref{eq:price-mass} and \eqref{eq:dyadic-triangle-mass} satisfy
\[
 T_0\le O_r(M_x).
\]
The implicit constant depends only on $r$.
\end{proposition}

Section~\ref{sec:triangles} proves the proposition by repeatedly
reducing neighborhood weights while retaining a controlled fraction of
the triangle mass.  In view of \eqref{eq:triangle-reduction}, this
completes the dual inequality and Lemma~\ref{lem:adjustment}.  The
constants are chosen in this order: the sparse multiplier construction
determines $A$ and $C_1$; the triangle bound then determines $B$; and
finally $\Lambda_r$ is increased so that the
$O(\Lambda^2)\sum_iw_i y_i$ term is at most
$\Lambda^4\sum_iw_i y_i$.
